\section{Kernel estimates uniform in depth}\label{sec:uniform-kernel}
The preceding fixed-depth estimates alone do not justify increasing
$L$ with $d$. We now prove the required uniform replacements. Put
$N_*=L-1$, $\beta=2^{-N_*}$, and $a_*=\pi^2/8$; this $a_*$ is unrelated
to the output margin $a$.

\begin{lemma}[Uniform endpoint coefficient transfer]\label{lem:uniform-coeff}
There are absolute $C_0,C_1,c,C>0$ such that, for every dyadic
$0<\beta\le1/2$ and every odd $j\ge\exp(C_1/\beta)$,
\[
 c\beta j^{-1-\beta}\le [z^j]\Psi^{\circ N_*}(z)
       \le C\beta j^{-1-\beta}.
\]
More precisely, with $\delta_\beta=\exp(-C_0/\beta)$, the coefficient is
\[
 \frac{2a_*^{\beta-1}}{-\Gamma(-\beta)}
 \frac{\Gamma(j-\beta)}{\Gamma(j+1)}
 +O(\beta j^{-1-3\beta})+O(e^{-cj\delta_\beta}),
\]
with absolute constants.
\end{lemma}
\begin{proof}
Let $h(q)=1-\cos(\pi q/2)=a_*q^2+O(q^4)$. On a fixed small disk define
\[
 \phi(q)=a_*q\exp\left\{\sum_{t\ge0}2^{-t-1}
 \log\frac{h(h^{\circ t}(q))}{a_*h^{\circ t}(q)^2}\right\}.
\]
The quotient is analytic, even, and $1+O(q^2)$. Quadratic contraction
makes the series locally uniformly convergent, with its derivatives.
Reindexing gives $\phi(h(q))=\phi(q)^2$. The function $\phi$ is odd,
real analytic, and $\phi'(0)=a_*$. Its inverse $\psi$ is odd analytic
on another fixed disk and $\psi'(0)=a_*^{-1}$.
The endpoint map $g(q)=1-\Psi(1-q)$ is the inverse branch of $h$.
Consequently, on the principal slit sector,
\begin{equation}\label{eq:bottcher}
 1-\Psi^{\circ N_*}(1-q)=\psi(\phi(q)^\beta).
\end{equation}
Choose $C_0$ large enough that $|q|\le\delta_\beta$ puts
$\phi(q)^\beta$ in a fixed compact subdisk of the domain of $\psi$.
On the upper cut $z=1+t+i0$, $q=-t-i0$ and oddness gives
\[
 \Im\Psi^{\circ N_*}(1+t+i0)
 =a_*^{\beta-1}\sin(\pi\beta)t^\beta+O(\beta t^{3\beta}),
 \qquad 0<t\le\delta_\beta.
\]
The factor $\beta$ in the error is important. Expand the odd series
$\psi(w)=\sum_{j\text{ odd}}b_jw^j$ and use
$|\sin(j\pi\beta)|\le j\pi\beta$; the weighted absolute coefficient
sum is uniformly bounded on that compact disk. Replacing $\phi(-t)$
by $-a_*t$ contributes only $O(\beta^2t^{\beta+2})$.

For completeness, the coefficient contour can be chosen uniformly.
Every iterate is analytic on
$\Omega=\mathbb C\setminus(( -\infty,-1]\cup[1,\infty))$:
$\Psi$ maps this domain into $-1<\Re z<1$ and hence avoids the cuts.
There is an absolute $C_H$ such that
\[
 \sup_{z\in\Omega,\ |z|\le1+\delta}|\Psi(z)|
       \le1+C_H\sqrt\delta\qquad(0<\delta\le1).
\]
Indeed project to the unit disk, where positivity of the power series
bounds the modulus by one, and use the uniform one-half H\"older
estimate near the only boundary singularities $\pm1$.
Iterating this inequality from $\delta=\delta_\beta$ bounds all iterates
on $|z|=1+\delta_\beta$ by two after increasing $C_0$.
Deform Cauchy's coefficient contour to this circle and the two cut
segments. The outer contribution is $O(e^{-cj\delta_\beta})$;
oddness makes the cut contributions add at odd $j$. They equal
\[
 \frac2\pi a_*^{\beta-1}\sin(\pi\beta)
 \int_0^{\delta_\beta}t^\beta(1+t)^{-j-1}\,dt
 +O\left(\beta\int_0^{\delta_\beta}t^{3\beta}(1+t)^{-j-1}\,dt\right).
\]
Extend the leading integral to infinity. The beta integral and Euler's
reflection formula give the displayed expansion. Its relative error
is small for $j\ge\exp(C_1/\beta)$ with $C_1$ sufficiently large.
\end{proof}

\begin{lemma}[Uniform spectral and cusp bounds]\label{lem:uniform-spectrum}
For odd $r\ge\exp(C_1/\beta)$, $r\le D/8$, the bounds in
Lemma~\ref{lem:sharp} hold with absolute constants after inserting a
factor $\beta$ in the level mass and eigenvalue-to-tail ratio and a
factor $\beta^2$ in the second-moment upper bound. In particular,
\[
 \tau_r\asymp r^{-\beta},\qquad
 \lambda_{r+1}/\tau_r\asymp
       \frac{\beta}{r\binom D{r+2}}.
\]
For the matrix eigenvalue $v=\lambda_{\min}(K)$,
\[
 v\ge c(d\log d)^{-\beta},\qquad
 \lambda_0\ge\beta^{\log_2(\pi/2)}/D.
\]
Uniformly at every distance and every spectral cutoff $t\ge1$,
\[
 1-K(x,y)\le C(\dist(x,y)/d)^\beta,
 \qquad \tau_t\le C(t+1)^{-\beta}.
\]
\end{lemma}
\begin{proof}
In the positive walk expansion of the degree-$j$ eigenvalue, all walk
times below $j$ vanish. Lemma~\ref{lem:uniform-coeff} and the gamma
integral therefore give, uniformly for $j\ge r$ odd,
\[
 a_j\asymp\beta D^{-\beta}\int_0^\infty
 u^{-1-\beta}p(u)^j(1-p(u))^{D-j}\,du,
 \qquad p(u)=(1-e^{-2u})/2.
\]
The beta-integral proof of the level estimate is uniform for
$0<\beta\le1/2$. The integral away from $p=0$ has at most a factor
$1/\beta$, canceled by the prefactor, and its binomial term is
exponentially small. The resulting level mass is
$\binom Dj a_j\asymp\beta j^{-1-\beta}$ up to any fixed fraction below
$D/2$.
For the tail upper bound, integrate
$\beta u^{-1-\beta}\min(1,Du/r)$; its constant is uniform.
For the lower bound integrate over the whole interval $u\ge Cr/D$:
a $\operatorname{Bin}(D,p(u))$ variable exceeds $r$ with odd parity
with probability bounded below. This gives $cr^{-\beta}$; summing
only the first dyadic shell would incorrectly lose a factor $\beta$.
The paired ratio and second moment follow from the same binomial
ratios as in Lemma~\ref{lem:sharp}. The exponentially small far-degree
remainder is absorbed uniformly: the threshold
$r\ge\exp(C_1/\beta)$ implies $\beta^{-1}\le C\log r$,
whereas the remaining binomial gap is exponential in $d$.

To bound $v$, use paired monotonicity and an odd full-cube degree near
$D$. Restrict the integral to $u\ge\frac12\log(2D)$, where its binomial
monomial is at least $c2^{-D}$. The integral of
$\beta u^{-1-\beta}$ cancels the factor $1/\beta$. Multiplication by
$2^d$ gives the matrix bound. The constant mode is bounded below by
the linear Taylor term $(2/\pi)^{N_*}/D$.

For the cusp, \eqref{eq:bottcher} gives the claimed uniform bound when
$1-(1+x\cdot y)/D\le\delta_\beta$. Outside that range use the bound two
and $\delta_\beta^\beta=e^{-C_0}$. The full tail above
$\exp(C_1/\beta)$ was just proved; below that threshold the trivial
bound one is at most $e^{C_1}(t+1)^{-\beta}$.
\end{proof}

All subsequent uses of depth-dependent fixed-depth estimates can now
be made uniform. The first-shell Gram proof uses only the upper shell
mass, lower tail, paired ratios, and positive monotonicity. Radial
second moments use denominators $2-\beta$, bounded away from zero.
The angular-envelope Cauchy weights below use the fixed exponent
$1/4$, so no factor $1/\beta$ occurs in their sum. Representation,
Gaussian concentration, VC, and recurrence constants are independent
of depth. Every displayed positive-power error exponent is smallest
at $\beta=1/2$. The fine covariance witnesses use only uniform
\emph{upper} bounds, so their lower energy constants remain uniform.
The polynomial sparsity window absorbs all explicit logarithmic
factors. These observations will be used in the final assembly.
